\section{Build vs Buy vs Build-with：冰山以下的复杂性}

讲者提到一个高频客户问题：“为什么不自己做？”很多工程团队在白板上画出来的方案通常是：选模型、加向量库、接工具 API、上线。现实是，真正困难的部分在“冰山水下”：版本治理、发布回滚、策略审计、异常分流、法务约束、跨模型故障切换等。

\begin{table}[H]
\centering
\caption{Agent 平台“冰山模型”：显性组件与隐性能力}
\begin{tabularx}{\textwidth}{>{\raggedright\arraybackslash}p{3.4cm} >{\raggedright\arraybackslash}X}
\toprule
层级 & 典型内容 \\
\midrule
显性（容易被看到） & Model selection、RAG、tool calling、prompt engineering、UI channel integration \\
隐性（上线后最耗成本） & 策略版本管理、灰度发布、线上回滚、跨渠道状态一致性、合规规则、审计追踪、failover、异常输入防护、人工接管机制 \\
\bottomrule
\end{tabularx}
\end{table}

\begin{warningbox}{“能跑起来”不等于“能运营”}
许多 PoC 在小流量下看似稳定，一旦进入真实客户分布，长尾输入和系统抖动会迅速暴露架构短板。没有 release discipline 和 policy enforcement 的 Agent，很难跨过首个业务高峰。
\end{warningbox}

\begin{importantbox}{build-with 的中间路径}
讲者给出的实务路径不是二选一，而是 build-with：在平台提供的高层抽象上做业务配置，保留必要扩展能力，避免重复建设低层基础设施。
\end{importantbox}

\subsection{本章小结}

build/buy 的核心不是“谁写代码”，而是“谁承担长期复杂性”。能否系统化管理冰山水下能力，是企业 Agent 项目能否持续两年以上的分水岭。

